\documentclass{cvclass}

% -------------------- Sizes
\def\margin{0.4cm}
\def\headerHeight{5cm}
\def\spaceBetweenContent{0.2cm}
\def\linkSize{5cm}
\def\squareSize{\dimexpr\headerHeight-\margin*2\relax}
\def\triangleSize{\margin}

% Définition de \contentWidth et \blockWidth
\newlength{\contentWidth}
\setlength{\contentWidth}{\dimexpr\paperwidth - (\margin*2)\relax}
\newlength{\blockWidth}
\setlength{\blockWidth}{\dimexpr (\paperwidth - \margin*7)/2\relax}

% Définition de \topContent, \contentHeight, \textTop et \textHeight
\def\topContent{(\headerHeight + \spaceBetweenContent + \margin)}
\def\contentHeight{\dimexpr\paperheight - \topContent - \margin\relax}
\def\textTop{(\headerHeight+\margin+\spaceBetweenContent)}
\def\textHeight{(\contentHeight - \margin*2)}

% Espaces entre les lignes
\def\spaceSections{0.15cm}
% -------------------- Colors
\definecolor{primaryColor}{RGB}{0, 0, 0}
\definecolor{headerRed}{HTML}{691616}
\definecolor{TitleRed}{HTML}{FF001F}
\definecolor{jobs}{RGB}{105, 105, 105}



% -------------------- Content
\newcommand{\cvname}{Mélissa COLIN}
\newcommand{\cvdesc}{Élève-ingénieure en Informatique et Intelligence \\Artificielle à l’ENSEIRB-MATMECA.}
\newcommand{\github}{https://urlr.me/6Qy4Kc}
\newcommand{\medium}{https://urlr.me/GNY4s}
\newcommand{\email}{melissa.colin0@proton.me}
\newcommand{\location}{Bordeaux et périphérie}
\newcommand{\website}{https://melissacolin.ai}
\newcommand{\linkedin}{https://urlr.me/Bpqxd}
\newcommand{\phone}{+33 7 82 52 XX XX}
% \newcommand{\phone}{+33 7 82 52 XX XX}
\newcommand{\search}{Stage de recherche dans le domaine de l’Intelligence Artificielle, sur la période du 18 mai au 25 septembre 2026}
\newcommand{\profilepicture}{pdp.png}
\newcommand{\background}{background.jpg}



% -------------------- Document
\begin{document}

\header
% \@pdfcomment
\content


\begin{area-left}
  \section{Compétences}
    \subsection{Langages}
      Python, R, SQL
    
      \vspace{\spaceSections} 
      \subsection{Bibliothèques et frameworks}
      OpenCV, PyTorch, TensorFlow, NumPy, Pandas, scikit-learn, Grad-CAM, LIME, SHAP
    
      \vspace{\spaceSections} 
      \subsection{IA, data, et computer vision}
      Prétraitement des données, évaluation, explicabilité, Deep Learning, CNN, transformer, ViT, encoder-decoder, attention
    
      \vspace{\spaceSections} 
      \subsection{Logiciels et outils}
      Excel, Docker, Elasticsearch, LaTeX, Kubeflow, Kubernetes, Canva, Git/GitHub
    
      \vspace{\spaceSections} 
      \subsection{Langues}
      Français, anglais (B1), chinois (A1), espagnol (A2)
    
      \vspace{\spaceSections} 
      \subsection{Autre}
      Méthodes Agiles / SCRUM, RGPD, IA Act
    
      \vspace{\spaceSections} 
    \subsection{Soft Skills}
      Curiosité, organisation, leadership, collaboration, adaptabilité


\section{Publications}
  \textcolor{darkgray}{Colin, M. (2024) \textit{YOLOv8 Explained: Understanding Object Detection from Scratch}. Medium. \href{https://medium.com/@melissa.colin/yolov8-explained-understanding-object-detection-from-scratch-763479652312}{\textcolor{darkgray}{\faPaperclip}}}\\ \\
  \textbf{Colin, M.}, \& \textbf{Hraibi Kaadoud, I.} (2024). \textit{Performances et explicabilité de ViT et d’architectures CNN : une étude empirique utilisant LIME, SHAP et Grad-CAM}. RJCIA (PFIA 2024). \href{https://hal.science/hal-04641791v1}{\faPaperclip}

\section{Formations}
  \subsection{Diplôme d'ingénieur en informatique}
    \job{ENSEIRB-MATMECA (Bordeaux)}{09/2024 -- 09/2027}
    Parcours ingénieur-docteur
    % Spécialisation en Intelligence Artificielle prévue en 3ᵉ année.

    \vspace{\spaceSections}
    \subsection{Titre RNCP de niveau 6 : Développeuse IA}
    \job{EPSI (Bordeaux)}{09/2021 -- 09/2024}
    % VidAI : Outil web simplifiant l'édition vidéo grâce à une assistance IA.\\ \\
    % EcoSort : Application mobile pour le tri des déchets via reconnaissance d’image et géolocalisation, gagnante d'une compétition nationale d'innovation. \href{https://www.youtube.com/watch?v=GGjImtkW-us}{\faPaperclip}

    \vspace{\spaceSections}
    \subsection{BTS Services informatiques aux organisations}
    \job{EPSI (Bordeaux)}{09/2021 -- 09/2023}
    Option Solutions Logicielles et Applications Métiers


    % \vspace{\spaceSections}
    % \subsection{Baccalauréat Sciences et Technologies de l'Industrie et du Développement Durable}
    % \job{Lycée Les Iris (Lormont)}{09/2019 -- 09/2021}
    % Option SIN : Systèmes d'Information et Numérique 

\section{Centres d'intérêts}
  Audiovisuel : Montage vidéo, cadrage, prise et retouche photos\\
  Sport : Demi-fond et escalade
  \end{area-left}

\begin{area-right}
  \section{Expériences professionnelles}

    \subsection{Ambassadrice}
      \job{Capgemini}{07/2025 -- ..}
      Organisation de différents ateliers entre les étudiants de l'ENSEIRB-MATMECA et les équipes de Capgemini.
      \vspace{\spaceSections} 

    \subsection{Responsable Formations}
      \job{Eirb'IA (Talence)}{06/2025 -- ..}
      Organisation et animation de formations en IA
      \vspace{\spaceSections}
      
    \subsection{Vice Présidente}
      \job{Forum Ingenib (Talence)}{04/2025 -- ..}
        Coordination des équipes pour la mise en place du forum annuel de recrutement de l'ENSEIRB-MATMECA.
      \vspace{\spaceSections} 

      \subsection{Responsable Networking}
      \job{ai4industry (Talence)}{10/2024 -- ...}
      \begin{itemize}
        \item Planification, allocation des ressources, gestion budgétaire.
        \item Coordination des équipes et intervenants.
        \item Animation de l'événement (table ronde, discours d’ouverture/clôture).
      \end{itemize}
          \vspace{\spaceSections} 

    \subsection{Ingénieure R\&D - IA et sûreté (stage)}
      \job{Sector Group (Paris)}{06/2025 -- 08/2025}
      \begin{itemize}
        \item Sensibilisation et formation des managers à la mise en place de l'IA dans leurs équipes
        \item Conception et développement d'un système RAG (Retrieval-Augmented Generation) sécurisé pour l'analyse de connaissances internes.
        \item Intégration d'un pipeline d'analyse de la donnée selon son format.
        \item Étude de certifications de systèmes critiques (norme IEC 61508).
      \end{itemize}
      
    \subsection{Développeuse IA (alternance)}
      \job{Cali Intelligences (Talence)}{07/2023 -- 09/2024}
      \begin{leftarrowbox}
        \begin{itemize}
          \item Développement d'une nouvelle méthode IA de reconnaissance des comportements suspects.
          \item Amélioration des modèles (60\%  $\rightarrow$ 90\% de pertinence) et de leur compréhension (XAI : IA Explicable).
          \item Optimisation des hyperparamètres et entraînement de modèles de détection et de reconnaissance d’actions.
          \item Développement de pipelines d’entraînement (MLOps).
          \item Réduction du temps d’inférence des modèles de 70\%.
          \item Tri, analyse, statistiques et mise en place de processus de QA/QC pour une grande quantitée des données.
        \end{itemize}
      \end{leftarrowbox}
      \vspace{\spaceSections} 
      \job{Développeuse Python (stage)}{01/2023 -- 03/2023}
      \begin{dendleftarrowbox}
        \begin{itemize}
          \item Développement, tests et optimisation d’algorithmes en Pyhton.
          \item Entraînement de modèles de détection d’objets.
        \end{itemize}
      \end{dendleftarrowbox}\\

      \vspace{\spaceSections} 
    % \subsection{Permanence Contraception Masculine}
    % \job{Planning Familial / ENSC}{01/2022 -- 09/2024}

    \subsection{Chargée de projet web (stage)}
      \job{Eveho (Cenon)}{06/2022 -- 08/2022}

    % \subsection{Webmaster (stage)}
    %   \job{COACHINTERNET (Mérignac)}{05/2022 -- 06/2022}

    % \subsection{Professeure particulière de mathématiques}
    %   \job{Superprof}{12/2018 -- 12/2021}
        
\end{area-right}
\searchbar

\end{document}
